\begin{table}[t]
\centering
\caption{Oriented receptive fields against random initialization under R-STDP, by
decoder. The learned readout is the network's own decoder; the linear probe is an
external control fitted on the same features. The two agree in sign on every dataset but
not in magnitude, so any statement of this contrast has to name its decoder. Dispersion
is allowed to vary by dataset for the learned readout ($\phi \sim$ dataset); for the probe
the seed variance is not identifiable, so that model omits the seed intercept.}
\label{tab:phase2_prior}
\begin{tabular}{lrrrr}
\toprule
& \multicolumn{2}{c}{Learned readout} & \multicolumn{2}{c}{Linear probe} \\
\cmidrule(lr){2-3} \cmidrule(lr){4-5}
Dataset & $\Delta$ (pp) & $95\%$ CI & $\Delta$ (pp) & $95\%$ CI \\
\midrule
MNIST & $+4.39^{***}$ & $[+3.45,\, +5.33]$ & $+6.17^{***}$ & $[+5.81,\, +6.53]$ \\
FMNIST & $-3.90^{*}$ & $[-6.90,\, -0.86]$ & $-0.78^{**}$ & $[-1.32,\, -0.25]$ \\
KMNIST & $+3.51^{***}$ & $[+2.58,\, +4.43]$ & $+6.54^{***}$ & $[+5.96,\, +7.11]$ \\
notMNIST & $-7.57^{***}$ & $[-9.55,\, -5.55]$ & $-2.39^{***}$ & $[-3.28,\, -1.52]$ \\
\bottomrule
\end{tabular}

\footnotesize $^{*}p<0.05$, $^{**}p<0.01$, $^{***}p<0.001$.
\end{table}
